El presente análisis experimental demuestra empíricamente que la notación asintótica no es el único factor determinante en el rendimiento real de un algoritmo. Si bien las cotas teóricas orientan el comportamiento esperado, son las decisiones de implementación y su interacción a bajo nivel con la memoria del computador las que dictan el verdadero costo computacional.

En definitiva, la elección de un algoritmo en ingeniería de software debe ir más allá de su comportamiento asintótico e incluir obligatoriamente el tamaño real de la entrada, la distribución de los datos y, sobre todo, el costo físico que las estructuras temporales y los patrones de acceso generan sobre la memoria del computador.
